\documentclass{article}
\usepackage{../mathnotezh}

\title{海伦公式的推论}
\author{tbj}
\date{}

\begin{document}

\maketitle

已知三角形三边长求其面积,一般会使用海伦公式,可如果边长中带有根号,海伦公式的形式就复杂了,接下来介绍海伦公式的一个推论.

\begin{thm}
    设三角形$ABC$三边长为$a, b, c$,设$x, y, z$满足
    \[ \begin{cases}
        x + y = a^2 \\
        y + z = b^2 \\
        z + x = c^2
    \end{cases} \]

    则$Area = \frac{1}{2} \sqrt{xy + yz + zx}$.
\end{thm}

\begin{proof}
    由题意知
    \[ \begin{cases}
        x = \frac{1}{2}(c^2 + a^2 - b^2) \\
        y = \frac{1}{2}(a^2 + b^2 - c^2) \\
        z = \frac{1}{2}(b^2 + c^2 - a^2)
    \end{cases} \]

    因此
    \begin{align*}
        & \frac{1}{2} \sqrt{xy + yz + zx} \\
        & = \frac{1}{4} \sqrt{(c^2 + a^2 - b^2)(a^2 + b^2 - c^2) + (a^2 + b^2 - c^2)(b^2 + c^2 - a^2) +
            (b^2 + c^2 - a^2)(c^2 + a^2 - b^2)} \\
        & = \frac{1}{4} \sqrt{a^4 - (b - c)^2 + b^4 - (c - a)^2 + c^4 - (a - b)^2} \\
        & = \frac{1}{4} \sqrt{2a^2b^2 + 2b^2c^2 + 2c^2a^2 - a^4 - b^4 - c^4} \\
        & = \frac{1}{4} \sqrt{(a + b + c)(b + c - a)(c + a - b)(a + b - c)} \\
        & = \sqrt{p(p - a)(p - b)(p - c)} \\
        & = Area
    \end{align*}

    最后一步运用了海伦公式.
\end{proof}

\begin{rqq}
    由余弦定理知
    \[ \begin{cases}
        x = ac \cos B = \vec{BA} \cdot \vec{BC} \\
        y = ab \cos C = \vec{CA} \cdot \vec{CB} \\
        z = bc \cos A = \vec{AB} \cdot \vec{AC}
    \end{cases} \]

    于是知
    \[ Area = \frac{1}{2} \sqrt{(\vec{BA} \cdot \vec{BC})(\vec{CA} \cdot \vec{CB})
    + (\vec{CA} \cdot \vec{CB})(\vec{AB} \cdot \vec{AC}) + (\vec{AB} \cdot \vec{AC})(\vec{BA} \cdot \vec{BC})} \]
\end{rqq}

\end{document}